\chapter{MuSK Antibodies}

\section*{What Is This Test?}
MuSK antibodies target muscle-specific kinase, a protein involved in communication between nerves and skeletal muscle. They are associated with a subgroup of autoimmune myasthenia gravis, particularly when acetylcholine-receptor antibodies are absent.

\section*{Alternative Names}
Anti-MuSK; Muscle-Specific Kinase Antibody; MuSK Autoantibody.

\section*{Why This Test Is Ordered}
Testing is appropriate when fatigable weakness, ptosis, diplopia, bulbar weakness, or respiratory weakness suggests myasthenia gravis and the acetylcholine-receptor antibody test is negative or the clinical pattern is strongly suggestive.

\section*{How This Test Is Performed}
A blood sample is analyzed by a validated immunoassay. Results are interpreted with neurologic examination, electrodiagnostic testing, and other myasthenia antibodies when appropriate.

\section*{Preparation, Risks, and Limitations}
No fasting is required. Venipuncture is the usual risk. A negative result does not exclude seronegative myasthenia gravis, and low-level results require clinical correlation because assay performance differs between laboratories.

\section*{Understanding Results}
A positive MuSK antibody strongly supports autoimmune myasthenia gravis in a compatible presentation. MuSK-associated disease often has prominent bulbar, facial, neck, or respiratory involvement and requires specialist management.

\section*{References}
International consensus guidance on management of myasthenia gravis.

\clearpage
\section*{Understanding Results and Follow-up}
The laboratory report should identify the specimen, method, units, reference interval or decision limit, and whether the result is positive, negative, abnormal, or inconclusive. Reference intervals vary with age, sex, pregnancy, specimen type, assay, and laboratory; a result outside a range is not automatically a diagnosis, and a result within a range does not always exclude disease. Discuss unexpected results with the ordering clinician, who can decide whether repeat or confirmatory testing, treatment, or referral is appropriate. Do not change medicines or delay urgent care based on an isolated result.


\section*{Regulatory Status and Authorization Date}
Regulatory status applies to the specific assay, instrument, manufacturer, and intended use rather than to the test name alone. No single approval date can be assigned to this test category. For a commercial assay, verify the current product in the relevant regulator's database: the U.S. Food and Drug Administration (FDA) clearance, approval, or authorization; India's Central Drugs Standard Control Organisation (CDSCO) licence or permission; the European Union In Vitro Diagnostic Regulation (EU IVDR) CE certificate issued through the applicable conformity-assessment route; or the United Kingdom Medicines and Healthcare products Regulatory Agency (MHRA) registration and UKCA/CE status. Laboratory-developed tests may not have an individual FDA, CDSCO, EU IVDR, or MHRA approval date.
\clearpage
